%%%%%%%%%%%%%%%%%%%%%%%%%%%%%%%%%%%%%%%%%%%%%%%%%%%%%%%%%%%%
% 8.9 EXPERIMENTAL DESIGN
% The Composition of Advanced Mathematics
%%%%%%%%%%%%%%%%%%%%%%%%%%%%%%%%%%%%%%%%%%%%%%%%%%%%%%%%%%%%

\section{Experimental Design}
\label{sec:experimental-design}

\prerequisite{
Data collection and sampling, hypothesis testing, regression, analysis
of variance, probability, randomization, confidence intervals, effect
sizes, and basic optimization concepts.
}

\learningobjectives{
By the end of this section, the reader should be able to:
\begin{itemize}
    \item explain the purpose of experimental design;
    \item distinguish observational studies from designed experiments;
    \item identify experimental units, treatments, factors, and factor
          levels;
    \item explain the roles of randomization, replication, and blocking;
    \item distinguish random sampling from random assignment;
    \item understand completely randomized designs;
    \item understand randomized complete block designs;
    \item understand matched-pairs designs;
    \item recognize Latin square designs;
    \item understand factorial experiments;
    \item distinguish main effects and interaction effects;
    \item construct \(2^k\) factorial designs;
    \item understand coding factor levels as \(-1\) and \(+1\);
    \item calculate simple factorial effects;
    \item understand confounding in experimental designs;
    \item distinguish experimental confounding from observational
          confounding;
    \item understand fractional factorial designs conceptually;
    \item recognize aliasing;
    \item understand design resolution conceptually;
    \item understand nested and crossed factors;
    \item understand split-plot designs conceptually;
    \item understand repeated-measures designs;
    \item distinguish fixed and random factors;
    \item understand nuisance variables;
    \item distinguish replication from repeated measurement;
    \item recognize pseudoreplication;
    \item understand balance and orthogonality;
    \item understand treatment allocation ratios;
    \item understand power and sample-size considerations;
    \item understand precision gains from blocking;
    \item recognize crossover experiments;
    \item understand carryover effects conceptually;
    \item understand cluster-randomized experiments;
    \item understand interference and the stable-unit treatment-value
          assumption conceptually;
    \item understand adaptive and sequential design previews;
    \item identify common threats to experimental validity;
    \item connect experimental design with ANOVA, regression, causal
          inference, and optimization.
\end{itemize}
}

%%%%%%%%%%%%%%%%%%%%%%%%%%%%%%%%%%%%%%%%%%%%%%%%%%%%%%%%%%%%
\subsection{Why Design Experiments?}
\label{subsec:why-design-experiments}
%%%%%%%%%%%%%%%%%%%%%%%%%%%%%%%%%%%%%%%%%%%%%%%%%%%%%%%%%%%%

Statistical analysis cannot repair every weakness in how data were
generated.

A carefully designed experiment attempts to create data in a way that
makes scientific comparisons clear and statistically efficient.

The central idea is

\[
\boxed{
\text{design determines what conclusions the data can support}.
}
\]

Experimental design addresses questions such as:

\begin{itemize}
    \item Which treatments should be compared?
    \item Which experimental units should receive each treatment?
    \item How many observations are needed?
    \item Which nuisance variables should be blocked?
    \item Which factors should be varied together?
    \item How should randomization be implemented?
    \item What interactions should be estimable?
\end{itemize}

%%%%%%%%%%%%%%%%%%%%%%%%%%%%%%%%%%%%%%%%%%%%%%%%%%%%%%%%%%%%
\subsection{Designed Experiment}
\label{subsec:designed-experiment}
%%%%%%%%%%%%%%%%%%%%%%%%%%%%%%%%%%%%%%%%%%%%%%%%%%%%%%%%%%%%

\begin{definitionbox}{Designed Experiment}
A designed experiment deliberately assigns treatments or conditions to
experimental units according to a planned procedure and measures the
resulting responses.
\end{definitionbox}

The investigator controls at least part of the data-generating process.

This distinguishes an experiment from a purely observational study.

%%%%%%%%%%%%%%%%%%%%%%%%%%%%%%%%%%%%%%%%%%%%%%%%%%%%%%%%%%%%
\subsection{Experimental Unit}
\label{subsec:experimental-unit-design}
%%%%%%%%%%%%%%%%%%%%%%%%%%%%%%%%%%%%%%%%%%%%%%%%%%%%%%%%%%%%

\begin{definitionbox}{Experimental Unit}
An experimental unit is the smallest unit independently assigned to a
treatment condition.
\end{definitionbox}

Examples include:

\[
\boxed{
\text{a person},
}
\]

\[
\boxed{
\text{a plot of land},
}
\]

\[
\boxed{
\text{a machine},
}
\]

\[
\boxed{
\text{a classroom},
}
\]

or

\[
\boxed{
\text{a laboratory specimen}.
}
\]

Correctly identifying the experimental unit is essential for determining
the appropriate sample size and error term.

%%%%%%%%%%%%%%%%%%%%%%%%%%%%%%%%%%%%%%%%%%%%%%%%%%%%%%%%%%%%
\subsection{Observational Unit Versus Experimental Unit}
\label{subsec:observational-versus-experimental-unit}
%%%%%%%%%%%%%%%%%%%%%%%%%%%%%%%%%%%%%%%%%%%%%%%%%%%%%%%%%%%%

The observational unit is the unit on which measurements are recorded.

The experimental unit is the unit receiving independent treatment
assignment.

These may differ.

For example, suppose entire classrooms are randomly assigned to teaching
methods and individual student scores are measured.

Then:

\[
\boxed{
\text{experimental unit}
=
\text{classroom},
}
\]

while

\[
\boxed{
\text{observational unit}
=
\text{student}.
}
\]

Treating every student as an independently randomized experimental unit
would exaggerate the effective sample size.

%%%%%%%%%%%%%%%%%%%%%%%%%%%%%%%%%%%%%%%%%%%%%%%%%%%%%%%%%%%%
\subsection{Treatment}
\label{subsec:treatment-design}
%%%%%%%%%%%%%%%%%%%%%%%%%%%%%%%%%%%%%%%%%%%%%%%%%%%%%%%%%%%%

A treatment is a specific experimental condition imposed on an
experimental unit.

A treatment may involve one factor level or a combination of several
factor levels.

For example:

\[
\boxed{
\text{fertilizer A at high irrigation}
}
\]

is one treatment combination in a two-factor experiment.

%%%%%%%%%%%%%%%%%%%%%%%%%%%%%%%%%%%%%%%%%%%%%%%%%%%%%%%%%%%%
\subsection{Factor}
\label{subsec:factor-design}
%%%%%%%%%%%%%%%%%%%%%%%%%%%%%%%%%%%%%%%%%%%%%%%%%%%%%%%%%%%%

\begin{definitionbox}{Factor}
A factor is a controlled explanatory variable whose levels define
experimental conditions.
\end{definitionbox}

Examples include:

\[
\boxed{
\text{temperature},
}
\]

\[
\boxed{
\text{drug dose},
}
\]

\[
\boxed{
\text{machine speed},
}
\]

or

\[
\boxed{
\text{instructional method}.
}
\]

%%%%%%%%%%%%%%%%%%%%%%%%%%%%%%%%%%%%%%%%%%%%%%%%%%%%%%%%%%%%
\subsection{Factor Levels}
\label{subsec:factor-levels-design}
%%%%%%%%%%%%%%%%%%%%%%%%%%%%%%%%%%%%%%%%%%%%%%%%%%%%%%%%%%%%

The values or categories used for a factor are called its levels.

For example, a temperature factor may have levels

\[
\boxed{
20^\circ,
\quad
40^\circ,
\quad
60^\circ.
}
\]

A treatment factor may have levels

\[
\boxed{
\text{control},
\quad
\text{treatment A},
\quad
\text{treatment B}.
}
\]

%%%%%%%%%%%%%%%%%%%%%%%%%%%%%%%%%%%%%%%%%%%%%%%%%%%%%%%%%%%%
\subsection{Response Variable}
\label{subsec:response-variable-design}
%%%%%%%%%%%%%%%%%%%%%%%%%%%%%%%%%%%%%%%%%%%%%%%%%%%%%%%%%%%%

The response variable is the outcome measured after applying the
experimental treatment.

Examples include:

\[
\boxed{
\text{yield},
}
\]

\[
\boxed{
\text{reaction time},
}
\]

\[
\boxed{
\text{strength},
}
\]

\[
\boxed{
\text{blood pressure},
}
\]

or

\[
\boxed{
\text{test score}.
}
\]

A single experiment may contain several response variables.

%%%%%%%%%%%%%%%%%%%%%%%%%%%%%%%%%%%%%%%%%%%%%%%%%%%%%%%%%%%%
\subsection{The Three Classical Principles}
\label{subsec:classical-design-principles}
%%%%%%%%%%%%%%%%%%%%%%%%%%%%%%%%%%%%%%%%%%%%%%%%%%%%%%%%%%%%

Three fundamental principles of experimental design are:

\[
\boxed{
\text{randomization},
}
\]

\[
\boxed{
\text{replication},
}
\]

and

\[
\boxed{
\text{blocking}.
}
\]

These principles address different sources of uncertainty.

%%%%%%%%%%%%%%%%%%%%%%%%%%%%%%%%%%%%%%%%%%%%%%%%%%%%%%%%%%%%
\subsection{Randomization}
\label{subsec:randomization-design}
%%%%%%%%%%%%%%%%%%%%%%%%%%%%%%%%%%%%%%%%%%%%%%%%%%%%%%%%%%%%

Randomization uses chance to assign experimental units to treatments.

Its purpose is to reduce systematic allocation bias and break systematic
relationships between treatment assignment and uncontrolled variables.

Conceptually,

\[
\boxed{
\text{randomization}
\rightarrow
\text{comparable treatment groups in expectation}.
}
\]

%%%%%%%%%%%%%%%%%%%%%%%%%%%%%%%%%%%%%%%%%%%%%%%%%%%%%%%%%%%%
\subsection{Why Randomization Supports Causal Inference}
\label{subsec:randomization-causal}
%%%%%%%%%%%%%%%%%%%%%%%%%%%%%%%%%%%%%%%%%%%%%%%%%%%%%%%%%%%%

Without randomization, units receiving different treatments may differ
systematically before treatment begins.

Random assignment makes treatment allocation independent of pre-treatment
characteristics under the design.

Therefore treatment groups differ systematically because of assigned
treatment rather than because investigators deliberately selected
different types of units.

This is why randomized experiments provide strong evidence for causal
effects.

%%%%%%%%%%%%%%%%%%%%%%%%%%%%%%%%%%%%%%%%%%%%%%%%%%%%%%%%%%%%
\subsection{Random Sampling Versus Random Assignment}
\label{subsec:random-sampling-random-assignment-design}
%%%%%%%%%%%%%%%%%%%%%%%%%%%%%%%%%%%%%%%%%%%%%%%%%%%%%%%%%%%%

Random sampling and random assignment have different purposes.

Random sampling addresses

\[
\boxed{
\text{generalization from sample to population}.
}
\]

Random assignment addresses

\[
\boxed{
\text{causal comparison among treatments}.
}
\]

A study can have one without the other.

%%%%%%%%%%%%%%%%%%%%%%%%%%%%%%%%%%%%%%%%%%%%%%%%%%%%%%%%%%%%
\subsection{Replication}
\label{subsec:replication-design}
%%%%%%%%%%%%%%%%%%%%%%%%%%%%%%%%%%%%%%%%%%%%%%%%%%%%%%%%%%%%

\begin{definitionbox}{Replication}
Replication means independently applying each treatment to multiple
experimental units.
\end{definitionbox}

Replication allows treatment effects to be separated from natural
unit-to-unit variability.

It also provides information for estimating experimental error.

%%%%%%%%%%%%%%%%%%%%%%%%%%%%%%%%%%%%%%%%%%%%%%%%%%%%%%%%%%%%
\subsection{Why Replication Matters}
\label{subsec:why-replication-matters}
%%%%%%%%%%%%%%%%%%%%%%%%%%%%%%%%%%%%%%%%%%%%%%%%%%%%%%%%%%%%

Suppose treatment A is applied to one unit and treatment B to one unit.

If the responses differ, the difference could result from:

\[
\boxed{
\text{treatment},
}
\]

or simply

\[
\boxed{
\text{preexisting unit differences}.
}
\]

Replication provides repeated independent evidence within each treatment
condition.

%%%%%%%%%%%%%%%%%%%%%%%%%%%%%%%%%%%%%%%%%%%%%%%%%%%%%%%%%%%%
\subsection{Replication Versus Repeated Measurements}
\label{subsec:replication-versus-repeated-measurements}
%%%%%%%%%%%%%%%%%%%%%%%%%%%%%%%%%%%%%%%%%%%%%%%%%%%%%%%%%%%%

Repeated measurements on the same experimental unit do not necessarily
create additional independent replication.

Suppose one plant receives a treatment and is measured ten times.

There is still only

\[
\boxed{
1
}
\]

independently treated plant.

The ten measurements may improve measurement precision, but they do not
replace ten independently treated plants.

%%%%%%%%%%%%%%%%%%%%%%%%%%%%%%%%%%%%%%%%%%%%%%%%%%%%%%%%%%%%
\subsection{Pseudoreplication}
\label{subsec:pseudoreplication}
%%%%%%%%%%%%%%%%%%%%%%%%%%%%%%%%%%%%%%%%%%%%%%%%%%%%%%%%%%%%

\begin{definitionbox}{Pseudoreplication}
Pseudoreplication occurs when nonindependent measurements are analyzed as
though they were independent experimental replicates.
\end{definitionbox}

This often causes underestimated standard errors and exaggerated
statistical significance.

Examples include:

\[
\boxed{
\text{many measurements from one treated tank},
}
\]

\[
\boxed{
\text{many students from one randomized classroom},
}
\]

or

\[
\boxed{
\text{many subsamples from one experimental plot}.
}
\]

%%%%%%%%%%%%%%%%%%%%%%%%%%%%%%%%%%%%%%%%%%%%%%%%%%%%%%%%%%%%
\subsection{Blocking}
\label{subsec:blocking-design}
%%%%%%%%%%%%%%%%%%%%%%%%%%%%%%%%%%%%%%%%%%%%%%%%%%%%%%%%%%%%

\begin{definitionbox}{Blocking}
Blocking groups similar experimental units before treatment assignment
so that treatment comparisons are made within relatively homogeneous
groups.
\end{definitionbox}

Examples of blocking variables include:

\[
\boxed{
\text{age group},
}
\]

\[
\boxed{
\text{location},
}
\]

\[
\boxed{
\text{baseline performance},
}
\]

or

\[
\boxed{
\text{production batch}.
}
\]

%%%%%%%%%%%%%%%%%%%%%%%%%%%%%%%%%%%%%%%%%%%%%%%%%%%%%%%%%%%%
\subsection{Why Blocking Can Increase Precision}
\label{subsec:blocking-precision}
%%%%%%%%%%%%%%%%%%%%%%%%%%%%%%%%%%%%%%%%%%%%%%%%%%%%%%%%%%%%

Suppose an important nuisance variable explains substantial response
variation.

Without blocking, that variation contributes to the residual error.

With blocking, some of it is modeled separately.

Thus,

\[
\boxed{
\text{blocking}
\rightarrow
\text{smaller residual variation}
\rightarrow
\text{greater precision}.
}
\]

Blocking is especially useful when the blocking variable is strongly
related to the response.

%%%%%%%%%%%%%%%%%%%%%%%%%%%%%%%%%%%%%%%%%%%%%%%%%%%%%%%%%%%%
\subsection{Nuisance Variable}
\label{subsec:nuisance-variable}
%%%%%%%%%%%%%%%%%%%%%%%%%%%%%%%%%%%%%%%%%%%%%%%%%%%%%%%%%%%%

A nuisance variable affects the response but is not itself the primary
scientific focus.

Examples include:

\[
\boxed{
\text{day of measurement},
}
\]

\[
\boxed{
\text{operator},
}
\]

\[
\boxed{
\text{laboratory batch},
}
\]

or

\[
\boxed{
\text{location}.
}
\]

A nuisance variable can be:

\[
\boxed{
\text{controlled},
}
\]

\[
\boxed{
\text{randomized},
}
\]

\[
\boxed{
\text{blocked},
}
\]

or

\[
\boxed{
\text{modeled statistically}.
}
\]

%%%%%%%%%%%%%%%%%%%%%%%%%%%%%%%%%%%%%%%%%%%%%%%%%%%%%%%%%%%%
\subsection{Completely Randomized Design}
\label{subsec:completely-randomized-design}
%%%%%%%%%%%%%%%%%%%%%%%%%%%%%%%%%%%%%%%%%%%%%%%%%%%%%%%%%%%%

\begin{definitionbox}{Completely Randomized Design}
A completely randomized design assigns experimental units to treatments
entirely through a randomization procedure without additional blocking
restrictions.
\end{definitionbox}

Suppose there are \(N\) units and \(k\) treatments.

If treatment \(i\) receives \(n_i\) units, then

\[
\boxed{
\sum_{i=1}^{k}
n_i=N.
}
\]

%%%%%%%%%%%%%%%%%%%%%%%%%%%%%%%%%%%%%%%%%%%%%%%%%%%%%%%%%%%%
\subsection{Worked Example: Completely Randomized Design}
\label{subsec:worked-completely-randomized-design}
%%%%%%%%%%%%%%%%%%%%%%%%%%%%%%%%%%%%%%%%%%%%%%%%%%%%%%%%%%%%

\begin{workedexamplebox}{Three Treatments}

Suppose

\[
N=30
\]

experimental units are assigned equally to three treatments.

Then each treatment receives

\[
n_i=10.
\]

A valid design randomly partitions the \(30\) units into three groups of
size \(10\).

\begin{answerbox}
\[
\boxed{
10
\text{ independent experimental units per treatment}.
}
\]
\end{answerbox}

The resulting responses may be analyzed using one-way ANOVA if the
corresponding assumptions are appropriate.

\end{workedexamplebox}

%%%%%%%%%%%%%%%%%%%%%%%%%%%%%%%%%%%%%%%%%%%%%%%%%%%%%%%%%%%%
\subsection{Balanced Designs}
\label{subsec:balanced-design}
%%%%%%%%%%%%%%%%%%%%%%%%%%%%%%%%%%%%%%%%%%%%%%%%%%%%%%%%%%%%

A design is balanced when treatment combinations have equal or
systematically equal replication.

For a one-factor design,

\[
\boxed{
n_1=n_2=\cdots=n_k.
}
\]

Balanced designs often simplify:

\[
\boxed{
\text{analysis},
}
\]

\[
\boxed{
\text{interpretation},
}
\]

and

\[
\boxed{
\text{variance calculations}.
}
\]

%%%%%%%%%%%%%%%%%%%%%%%%%%%%%%%%%%%%%%%%%%%%%%%%%%%%%%%%%%%%
\subsection{Advantages of Balance}
\label{subsec:advantages-balance}
%%%%%%%%%%%%%%%%%%%%%%%%%%%%%%%%%%%%%%%%%%%%%%%%%%%%%%%%%%%%

Balanced designs often provide:

\[
\boxed{
\text{orthogonal treatment comparisons},
}
\]

\[
\boxed{
\text{similar precision across treatments},
}
\]

and

\[
\boxed{
\text{simpler ANOVA decompositions}.
}
\]

Balance is desirable but not always possible because of cost,
availability, attrition, or ethical constraints.

%%%%%%%%%%%%%%%%%%%%%%%%%%%%%%%%%%%%%%%%%%%%%%%%%%%%%%%%%%%%
\subsection{Unequal Allocation}
\label{subsec:unequal-treatment-allocation}
%%%%%%%%%%%%%%%%%%%%%%%%%%%%%%%%%%%%%%%%%%%%%%%%%%%%%%%%%%%%

Equal group sizes are not always optimal.

One treatment may be:

\[
\boxed{
\text{more expensive},
}
\]

\[
\boxed{
\text{scarce},
}
\]

or

\[
\boxed{
\text{more variable}.
}
\]

Unequal allocation can sometimes improve efficiency or reduce cost.

The design should account for the objective being optimized.

%%%%%%%%%%%%%%%%%%%%%%%%%%%%%%%%%%%%%%%%%%%%%%%%%%%%%%%%%%%%
\subsection{Randomized Complete Block Design}
\label{subsec:randomized-complete-block-design-detail}
%%%%%%%%%%%%%%%%%%%%%%%%%%%%%%%%%%%%%%%%%%%%%%%%%%%%%%%%%%%%

A randomized complete block design, abbreviated RCBD, divides units into
blocks and applies every treatment within every block.

If there are \(t\) treatments and \(b\) blocks, a standard model is

\[
\boxed{
Y_{ij}
=
\mu
+
\tau_i
+
\beta_j
+
\varepsilon_{ij}.
}
\]

Here:

\[
\tau_i
=
\text{treatment effect},
\]

and

\[
\beta_j
=
\text{block effect}.
\]

%%%%%%%%%%%%%%%%%%%%%%%%%%%%%%%%%%%%%%%%%%%%%%%%%%%%%%%%%%%%
\subsection{Randomization Within Blocks}
\label{subsec:randomization-within-blocks}
%%%%%%%%%%%%%%%%%%%%%%%%%%%%%%%%%%%%%%%%%%%%%%%%%%%%%%%%%%%%

Within each block, treatments are assigned randomly.

Thus units are compared under similar nuisance conditions.

For example, if soil quality varies substantially across a field, nearby
plots might form blocks.

Treatments are randomized among plots within each block.

%%%%%%%%%%%%%%%%%%%%%%%%%%%%%%%%%%%%%%%%%%%%%%%%%%%%%%%%%%%%
\subsection{Worked Example: Blocking}
\label{subsec:worked-blocking-design}
%%%%%%%%%%%%%%%%%%%%%%%%%%%%%%%%%%%%%%%%%%%%%%%%%%%%%%%%%%%%

\begin{workedexamplebox}{Four Fertilizers Across Five Fields}

Suppose four fertilizers are compared across five regions of a field.

Each region forms a block containing four plots.

Within every block, randomly assign the four fertilizers so each appears
once.

The design therefore contains

\[
5\times4=20
\]

experimental plots.

\begin{answerbox}
\[
\boxed{
\text{Treatment comparisons are made after accounting for block
differences.}
}
\]
\end{answerbox}

\end{workedexamplebox}

%%%%%%%%%%%%%%%%%%%%%%%%%%%%%%%%%%%%%%%%%%%%%%%%%%%%%%%%%%%%
\subsection{Matched-Pairs Design}
\label{subsec:matched-pairs-design-experimental}
%%%%%%%%%%%%%%%%%%%%%%%%%%%%%%%%%%%%%%%%%%%%%%%%%%%%%%%%%%%%

A matched-pairs design is a block design in which each block contains
two closely related experimental units or two conditions on the same
unit.

Examples include:

\[
\boxed{
\text{twins},
}
\]

\[
\boxed{
\text{matched subjects},
}
\]

or

\[
\boxed{
\text{before-and-after measurements}.
}
\]

%%%%%%%%%%%%%%%%%%%%%%%%%%%%%%%%%%%%%%%%%%%%%%%%%%%%%%%%%%%%
\subsection{Paired-Difference Model}
\label{subsec:paired-difference-design}
%%%%%%%%%%%%%%%%%%%%%%%%%%%%%%%%%%%%%%%%%%%%%%%%%%%%%%%%%%%%

For pair \(i\), define

\[
\boxed{
D_i
=
Y_{iA}-Y_{iB}.
}
\]

Then the treatment comparison is reduced to analysis of

\[
D_1,\ldots,D_n.
\]

The paired mean difference is

\[
\boxed{
\overline{D}.
}
\]

This can substantially reduce variance when paired responses are highly
positively correlated.

%%%%%%%%%%%%%%%%%%%%%%%%%%%%%%%%%%%%%%%%%%%%%%%%%%%%%%%%%%%%
\subsection{Latin Square Design}
\label{subsec:latin-square-design}
%%%%%%%%%%%%%%%%%%%%%%%%%%%%%%%%%%%%%%%%%%%%%%%%%%%%%%%%%%%%

A Latin square design controls two blocking variables while studying one
treatment factor.

For \(t\) treatments, construct a

\[
t\times t
\]

array in which every treatment appears exactly once in each row and once
in each column.

Rows represent one blocking factor.

Columns represent another.

%%%%%%%%%%%%%%%%%%%%%%%%%%%%%%%%%%%%%%%%%%%%%%%%%%%%%%%%%%%%
\subsection{Example Latin Square}
\label{subsec:example-latin-square}
%%%%%%%%%%%%%%%%%%%%%%%%%%%%%%%%%%%%%%%%%%%%%%%%%%%%%%%%%%%%

A \(4\times4\) Latin square may be

\[
\begin{array}{c|cccc}
&
1&2&3&4\\
\hline
1&A&B&C&D\\
2&B&C&D&A\\
3&C&D&A&B\\
4&D&A&B&C
\end{array}
\]

Each treatment appears once in every row and every column.

%%%%%%%%%%%%%%%%%%%%%%%%%%%%%%%%%%%%%%%%%%%%%%%%%%%%%%%%%%%%
\subsection{Latin Square Model}
\label{subsec:latin-square-model}
%%%%%%%%%%%%%%%%%%%%%%%%%%%%%%%%%%%%%%%%%%%%%%%%%%%%%%%%%%%%

A classical additive model is

\[
\boxed{
Y_{ijk}
=
\mu
+
\tau_i
+
\rho_j
+
\kappa_k
+
\varepsilon_{ijk},
}
\]

where:

\[
\tau_i
=
\text{treatment effect},
\]

\[
\rho_j
=
\text{row effect},
\]

and

\[
\kappa_k
=
\text{column effect}.
\]

Traditional Latin-square analysis assumes interactions among these
effects are negligible or not separately estimable.

%%%%%%%%%%%%%%%%%%%%%%%%%%%%%%%%%%%%%%%%%%%%%%%%%%%%%%%%%%%%
\subsection{Factorial Experiments}
\label{subsec:factorial-experiments}
%%%%%%%%%%%%%%%%%%%%%%%%%%%%%%%%%%%%%%%%%%%%%%%%%%%%%%%%%%%%

A factorial experiment studies multiple factors simultaneously.

If factor \(A\) has \(a\) levels and factor \(B\) has \(b\) levels, a
full factorial design includes every combination.

The number of treatment combinations is

\[
\boxed{
ab.
}
\]

With three factors having \(a\), \(b\), and \(c\) levels, there are

\[
\boxed{
abc
}
\]

combinations.

%%%%%%%%%%%%%%%%%%%%%%%%%%%%%%%%%%%%%%%%%%%%%%%%%%%%%%%%%%%%
\subsection{Why Factorial Designs Are Powerful}
\label{subsec:why-factorial-designs}
%%%%%%%%%%%%%%%%%%%%%%%%%%%%%%%%%%%%%%%%%%%%%%%%%%%%%%%%%%%%

Factorial designs estimate:

\[
\boxed{
\text{main effects}
}
\]

and

\[
\boxed{
\text{interactions}
}
\]

using the same experiment.

Studying factors separately cannot reveal interaction as efficiently.

Factorial experiments therefore often provide more information per
experimental unit.

%%%%%%%%%%%%%%%%%%%%%%%%%%%%%%%%%%%%%%%%%%%%%%%%%%%%%%%%%%%%
\subsection{Main Effect}
\label{subsec:main-effect-design}
%%%%%%%%%%%%%%%%%%%%%%%%%%%%%%%%%%%%%%%%%%%%%%%%%%%%%%%%%%%%

A main effect describes the average change in response caused by changing
one factor while averaging over levels of the other factors.

For two-level factor \(A\),

\[
\boxed{
\text{effect of }A
=
\overline{Y}_{A+}
-
\overline{Y}_{A-}.
}
\]

%%%%%%%%%%%%%%%%%%%%%%%%%%%%%%%%%%%%%%%%%%%%%%%%%%%%%%%%%%%%
\subsection{Interaction}
\label{subsec:interaction-design}
%%%%%%%%%%%%%%%%%%%%%%%%%%%%%%%%%%%%%%%%%%%%%%%%%%%%%%%%%%%%

Interaction occurs when the effect of one factor depends on the level of
another.

For factors \(A\) and \(B\),

\[
\boxed{
AB
\text{ interaction}
}
\]

means the difference between high and low \(A\) is not constant across
levels of \(B\).

Interaction is one of the most important reasons to use factorial
designs.

%%%%%%%%%%%%%%%%%%%%%%%%%%%%%%%%%%%%%%%%%%%%%%%%%%%%%%%%%%%%
\subsection{Two-Level Factorial Designs}
\label{subsec:two-level-factorial-designs}
%%%%%%%%%%%%%%%%%%%%%%%%%%%%%%%%%%%%%%%%%%%%%%%%%%%%%%%%%%%%

Suppose each of \(k\) factors has two levels.

A complete factorial experiment contains

\[
\boxed{
2^k
}
\]

treatment combinations.

Examples include:

\[
2^2=4,
\]

\[
2^3=8,
\]

and

\[
2^4=16.
\]

%%%%%%%%%%%%%%%%%%%%%%%%%%%%%%%%%%%%%%%%%%%%%%%%%%%%%%%%%%%%
\subsection{Coding Two-Level Factors}
\label{subsec:coding-two-level-factors}
%%%%%%%%%%%%%%%%%%%%%%%%%%%%%%%%%%%%%%%%%%%%%%%%%%%%%%%%%%%%

Two-level factors are often coded as

\[
\boxed{
-1
}
\]

for the low level and

\[
\boxed{
+1
}
\]

for the high level.

This coding simplifies estimation of main effects and interactions.

For a \(2^2\) design:

\[
\begin{array}{c|cc}
\text{Run}
&
A
&
B
\\
\hline
1&-1&-1\\
2&+1&-1\\
3&-1&+1\\
4&+1&+1
\end{array}
\]

%%%%%%%%%%%%%%%%%%%%%%%%%%%%%%%%%%%%%%%%%%%%%%%%%%%%%%%%%%%%
\subsection{Interaction Column}
\label{subsec:factorial-interaction-column}
%%%%%%%%%%%%%%%%%%%%%%%%%%%%%%%%%%%%%%%%%%%%%%%%%%%%%%%%%%%%

With \(-1,+1\) coding, the \(AB\) interaction column is formed by
multiplying the \(A\) and \(B\) columns.

Thus:

\[
\boxed{
AB=A\times B.
}
\]

For the \(2^2\) design:

\[
\begin{array}{c|ccc}
\text{Run}
&A&B&AB\\
\hline
1&-1&-1&+1\\
2&+1&-1&-1\\
3&-1&+1&-1\\
4&+1&+1&+1
\end{array}
\]

%%%%%%%%%%%%%%%%%%%%%%%%%%%%%%%%%%%%%%%%%%%%%%%%%%%%%%%%%%%%
\subsection{Worked Example: Main Effect in a \(2^2\) Design}
\label{subsec:worked-main-effect-2x2}
%%%%%%%%%%%%%%%%%%%%%%%%%%%%%%%%%%%%%%%%%%%%%%%%%%%%%%%%%%%%

\begin{workedexamplebox}{Computing the Effect of Factor \(A\)}

Suppose the four treatment-combination means are

\[
\begin{array}{c|cc}
&
B-
&
B+
\\
\hline
A-
&
10
&
14
\\
A+
&
18
&
30
\end{array}
\]

The mean response at high \(A\) is

\[
\overline{Y}_{A+}
=
\frac{
18+30
}{2}
=
24.
\]

At low \(A\),

\[
\overline{Y}_{A-}
=
\frac{
10+14
}{2}
=
12.
\]

Therefore,

\begin{answerbox}
\[
\boxed{
\text{Effect of }A
=
24-12
=
12.
}
\]
\end{answerbox}

\end{workedexamplebox}

%%%%%%%%%%%%%%%%%%%%%%%%%%%%%%%%%%%%%%%%%%%%%%%%%%%%%%%%%%%%
\subsection{Worked Example: Interaction Effect}
\label{subsec:worked-interaction-effect-2x2}
%%%%%%%%%%%%%%%%%%%%%%%%%%%%%%%%%%%%%%%%%%%%%%%%%%%%%%%%%%%%

\begin{workedexamplebox}{Computing an \(AB\) Interaction}

Using the same treatment means,

\[
\begin{array}{c|cc}
&
B-
&
B+
\\
\hline
A-
&
10
&
14
\\
A+
&
18
&
30
\end{array}
\]

the effect of \(A\) when \(B-\) is

\[
18-10=8.
\]

The effect of \(A\) when \(B+\) is

\[
30-14=16.
\]

A difference-of-differences interaction measure is

\[
16-8=8.
\]

Under the common factorial-effect convention that averages this
difference across the two levels,

\begin{answerbox}
\[
\boxed{
AB\text{ effect}=4.
}
\]
\end{answerbox}

The precise numerical scaling depends on the chosen factorial-effect
convention, so the convention should be stated.

\end{workedexamplebox}

%%%%%%%%%%%%%%%%%%%%%%%%%%%%%%%%%%%%%%%%%%%%%%%%%%%%%%%%%%%%
\subsection{Regression Representation of a \(2^2\) Factorial}
\label{subsec:regression-2x2-factorial}
%%%%%%%%%%%%%%%%%%%%%%%%%%%%%%%%%%%%%%%%%%%%%%%%%%%%%%%%%%%%

Using coded factors

\[
A,B\in\{-1,+1\},
\]

a model is

\[
\boxed{
Y
=
\beta_0
+
\beta_AA
+
\beta_BB
+
\beta_{AB}AB
+
\varepsilon.
}
\]

Under standard effect coding,

\[
2\beta_A
\]

corresponds to the main effect of \(A\),

\[
2\beta_B
\]

to the main effect of \(B\), and

\[
2\beta_{AB}
\]

to the interaction effect under the corresponding convention.

%%%%%%%%%%%%%%%%%%%%%%%%%%%%%%%%%%%%%%%%%%%%%%%%%%%%%%%%%%%%
\subsection{Three-Factor Designs}
\label{subsec:three-factor-designs}
%%%%%%%%%%%%%%%%%%%%%%%%%%%%%%%%%%%%%%%%%%%%%%%%%%%%%%%%%%%%

A \(2^3\) design contains factors

\[
A,
\qquad
B,
\qquad
C.
\]

The model may include

\[
\boxed{
A,\ B,\ C,
}
\]

two-factor interactions

\[
\boxed{
AB,\ AC,\ BC,
}
\]

and the three-factor interaction

\[
\boxed{
ABC.
}
\]

Higher-order interactions describe how lower-order effects depend on
additional factors.

%%%%%%%%%%%%%%%%%%%%%%%%%%%%%%%%%%%%%%%%%%%%%%%%%%%%%%%%%%%%
\subsection{Hierarchy Principle}
\label{subsec:hierarchy-principle}
%%%%%%%%%%%%%%%%%%%%%%%%%%%%%%%%%%%%%%%%%%%%%%%%%%%%%%%%%%%%

A common modeling principle is that if a high-order interaction is
included, the lower-order terms that compose it should generally remain
in the model.

For example, if

\[
AB
\]

is included, one usually retains

\[
A
\]

and

\[
B.
\]

This is called model hierarchy.

%%%%%%%%%%%%%%%%%%%%%%%%%%%%%%%%%%%%%%%%%%%%%%%%%%%%%%%%%%%%
\subsection{Effect Sparsity Principle}
\label{subsec:effect-sparsity}
%%%%%%%%%%%%%%%%%%%%%%%%%%%%%%%%%%%%%%%%%%%%%%%%%%%%%%%%%%%%

In many industrial experiments, only a small fraction of possible
effects are practically important.

This motivates the effect sparsity principle:

\[
\boxed{
\text{most high-order factorial effects are often negligible}.
}
\]

This principle motivates fractional factorial designs.

It is a modeling assumption, not a mathematical law.

%%%%%%%%%%%%%%%%%%%%%%%%%%%%%%%%%%%%%%%%%%%%%%%%%%%%%%%%%%%%
\subsection{Fractional Factorial Designs}
\label{subsec:fractional-factorial-designs}
%%%%%%%%%%%%%%%%%%%%%%%%%%%%%%%%%%%%%%%%%%%%%%%%%%%%%%%%%%%%

A full \(2^k\) factorial design can become large when \(k\) is large.

A fractional factorial design uses only a fraction of the full set of
treatment combinations.

Examples include

\[
\boxed{
2^{k-1},
}
\]

\[
\boxed{
2^{k-2},
}
\]

or smaller fractions.

The cost reduction comes at the price of confounding some effects with
others.

%%%%%%%%%%%%%%%%%%%%%%%%%%%%%%%%%%%%%%%%%%%%%%%%%%%%%%%%%%%%
\subsection{Aliasing}
\label{subsec:factorial-aliasing}
%%%%%%%%%%%%%%%%%%%%%%%%%%%%%%%%%%%%%%%%%%%%%%%%%%%%%%%%%%%%

In a fractional factorial design, two or more effects may be impossible
to distinguish from one another.

These effects are said to be aliased.

For example, a design may imply an alias relation such as

\[
\boxed{
A
\equiv
BCD.
}
\]

Then the observed contrast associated with \(A\) also contains the
\(BCD\) interaction.

%%%%%%%%%%%%%%%%%%%%%%%%%%%%%%%%%%%%%%%%%%%%%%%%%%%%%%%%%%%%
\subsection{Defining Relation Preview}
\label{subsec:defining-relation-preview}
%%%%%%%%%%%%%%%%%%%%%%%%%%%%%%%%%%%%%%%%%%%%%%%%%%%%%%%%%%%%

Fractional factorial designs are often generated using relations such as

\[
\boxed{
I=ABC.
}
\]

This relation identifies which columns are intentionally confounded.

Multiplying both sides by \(A\) gives

\[
A=BC.
\]

Thus \(A\) is aliased with \(BC\).

Full defining-relation theory is developed in more advanced design of
experiments.

%%%%%%%%%%%%%%%%%%%%%%%%%%%%%%%%%%%%%%%%%%%%%%%%%%%%%%%%%%%%
\subsection{Design Resolution Preview}
\label{subsec:design-resolution-preview}
%%%%%%%%%%%%%%%%%%%%%%%%%%%%%%%%%%%%%%%%%%%%%%%%%%%%%%%%%%%%

The resolution of a two-level fractional factorial design indicates the
minimum order of effects that are aliased with one another.

Conceptually:

\[
\boxed{
\text{Resolution III}
}
\]

can alias main effects with two-factor interactions,

\[
\boxed{
\text{Resolution IV}
}
\]

typically separates main effects from two-factor interactions but may
alias two-factor interactions with one another, and

\[
\boxed{
\text{Resolution V}
}
\]

typically separates main effects and two-factor interactions from each
other more effectively.

%%%%%%%%%%%%%%%%%%%%%%%%%%%%%%%%%%%%%%%%%%%%%%%%%%%%%%%%%%%%
\subsection{Experimental Confounding}
\label{subsec:experimental-confounding}
%%%%%%%%%%%%%%%%%%%%%%%%%%%%%%%%%%%%%%%%%%%%%%%%%%%%%%%%%%%%

In design of experiments, confounding can have a deliberate technical
meaning.

Two effects are confounded if the design does not provide independent
information to estimate them separately.

This differs from observational confounding, where a third variable
creates ambiguity in interpreting an exposure-outcome relationship.

Thus the same word appears in two related but distinct contexts.

%%%%%%%%%%%%%%%%%%%%%%%%%%%%%%%%%%%%%%%%%%%%%%%%%%%%%%%%%%%%
\subsection{Orthogonality}
\label{subsec:orthogonality-design}
%%%%%%%%%%%%%%%%%%%%%%%%%%%%%%%%%%%%%%%%%%%%%%%%%%%%%%%%%%%%

A design is orthogonal when effect-estimation columns are orthogonal.

For coded design columns \(x_j\) and \(x_k\),

\[
\boxed{
x_j^Tx_k=0
}
\]

for distinct effects.

Orthogonality leads to independent or uncorrelated effect estimates under
the classical model and simplifies sums-of-squares calculations.

%%%%%%%%%%%%%%%%%%%%%%%%%%%%%%%%%%%%%%%%%%%%%%%%%%%%%%%%%%%%
\subsection{Why Orthogonality Is Valuable}
\label{subsec:why-orthogonality}
%%%%%%%%%%%%%%%%%%%%%%%%%%%%%%%%%%%%%%%%%%%%%%%%%%%%%%%%%%%%

When effects are orthogonal:

\[
\boxed{
\text{one effect can be estimated without changing another},
}
\]

and

\[
\boxed{
\text{sums of squares partition cleanly}.
}
\]

Balanced factorial designs frequently possess strong orthogonality
properties.

%%%%%%%%%%%%%%%%%%%%%%%%%%%%%%%%%%%%%%%%%%%%%%%%%%%%%%%%%%%%
\subsection{Randomization Restrictions}
\label{subsec:randomization-restrictions}
%%%%%%%%%%%%%%%%%%%%%%%%%%%%%%%%%%%%%%%%%%%%%%%%%%%%%%%%%%%%

Sometimes complete randomization is impractical.

For example, one factor may be difficult or expensive to change.

This motivates restricted-randomization designs such as:

\[
\boxed{
\text{blocks},
}
\]

\[
\boxed{
\text{split plots},
}
\]

and

\[
\boxed{
\text{clusters}.
}
\]

Restricted randomization changes the correct error structure.

%%%%%%%%%%%%%%%%%%%%%%%%%%%%%%%%%%%%%%%%%%%%%%%%%%%%%%%%%%%%
\subsection{Split-Plot Design}
\label{subsec:split-plot-design}
%%%%%%%%%%%%%%%%%%%%%%%%%%%%%%%%%%%%%%%%%%%%%%%%%%%%%%%%%%%%

A split-plot design arises when one factor is difficult to randomize at
the individual-unit level.

The hard-to-change factor is assigned to larger units called whole
plots.

Another factor is randomized within each whole plot to subplots.

Thus there are two randomization levels.

%%%%%%%%%%%%%%%%%%%%%%%%%%%%%%%%%%%%%%%%%%%%%%%%%%%%%%%%%%%%
\subsection{Why Split-Plot Analyses Need Multiple Error Terms}
\label{subsec:split-plot-error-terms}
%%%%%%%%%%%%%%%%%%%%%%%%%%%%%%%%%%%%%%%%%%%%%%%%%%%%%%%%%%%%

Because whole-plot treatments are randomized at one level and subplot
treatments at another, the precision for different effects is not the
same.

Whole-plot effects are estimated using between-whole-plot variability.

Subplot effects use within-whole-plot variability.

Therefore a single undifferentiated residual variance is generally
incorrect.

%%%%%%%%%%%%%%%%%%%%%%%%%%%%%%%%%%%%%%%%%%%%%%%%%%%%%%%%%%%%
\subsection{Nested Designs}
\label{subsec:nested-design}
%%%%%%%%%%%%%%%%%%%%%%%%%%%%%%%%%%%%%%%%%%%%%%%%%%%%%%%%%%%%

Factor \(B\) is nested within factor \(A\) when each level of \(B\)
belongs to exactly one level of \(A\).

For example,

\[
\boxed{
\text{students nested within classrooms},
}
\]

or

\[
\boxed{
\text{machines nested within factories}.
}
\]

Nested structures naturally arise in hierarchical experiments.

%%%%%%%%%%%%%%%%%%%%%%%%%%%%%%%%%%%%%%%%%%%%%%%%%%%%%%%%%%%%
\subsection{Crossed Factors}
\label{subsec:crossed-factors-design}
%%%%%%%%%%%%%%%%%%%%%%%%%%%%%%%%%%%%%%%%%%%%%%%%%%%%%%%%%%%%

Factors are crossed when each level of one factor appears with every
level of the other factor.

For factors \(A\) and \(B\),

\[
\boxed{
A\times B
}
\]

indicates a crossed factorial structure.

Crossing allows direct estimation of interactions between the factors.

%%%%%%%%%%%%%%%%%%%%%%%%%%%%%%%%%%%%%%%%%%%%%%%%%%%%%%%%%%%%
\subsection{Fixed Factors}
\label{subsec:fixed-factors-design}
%%%%%%%%%%%%%%%%%%%%%%%%%%%%%%%%%%%%%%%%%%%%%%%%%%%%%%%%%%%%

A factor is fixed when the specific levels used in the experiment are
the levels of scientific interest.

For example, comparing three named machines or three specific drugs often
uses fixed effects.

Inference concerns those particular levels.

%%%%%%%%%%%%%%%%%%%%%%%%%%%%%%%%%%%%%%%%%%%%%%%%%%%%%%%%%%%%
\subsection{Random Factors}
\label{subsec:random-factors-design}
%%%%%%%%%%%%%%%%%%%%%%%%%%%%%%%%%%%%%%%%%%%%%%%%%%%%%%%%%%%%

A factor may be random when observed levels are regarded as sampled from
a larger population of possible levels.

For example, a study may randomly sample production batches.

Then the batch effect may be modeled as

\[
\boxed{
B_j
\sim
N(0,\sigma_B^2).
}
\]

The main interest may be the variance component

\[
\sigma_B^2.
\]

%%%%%%%%%%%%%%%%%%%%%%%%%%%%%%%%%%%%%%%%%%%%%%%%%%%%%%%%%%%%
\subsection{Mixed-Effects Designs}
\label{subsec:mixed-effects-designs}
%%%%%%%%%%%%%%%%%%%%%%%%%%%%%%%%%%%%%%%%%%%%%%%%%%%%%%%%%%%%

A mixed-effects design contains both fixed and random factors.

For example,

\[
\boxed{
\text{treatment}
=
\text{fixed},
}
\]

while

\[
\boxed{
\text{site}
=
\text{random}.
}
\]

Mixed models are often appropriate for multisite, repeated-measure, and
clustered experiments.

%%%%%%%%%%%%%%%%%%%%%%%%%%%%%%%%%%%%%%%%%%%%%%%%%%%%%%%%%%%%
\subsection{Repeated-Measures Designs}
\label{subsec:repeated-measures-design}
%%%%%%%%%%%%%%%%%%%%%%%%%%%%%%%%%%%%%%%%%%%%%%%%%%%%%%%%%%%%

In a repeated-measures design, the same experimental units are observed
under several conditions or times.

This creates within-unit correlation.

Advantages include:

\[
\boxed{
\text{control of stable subject differences},
}
\]

and

\[
\boxed{
\text{greater precision}.
}
\]

However, the dependence structure must be modeled.

%%%%%%%%%%%%%%%%%%%%%%%%%%%%%%%%%%%%%%%%%%%%%%%%%%%%%%%%%%%%
\subsection{Order Effects}
\label{subsec:order-effects}
%%%%%%%%%%%%%%%%%%%%%%%%%%%%%%%%%%%%%%%%%%%%%%%%%%%%%%%%%%%%

When each unit receives several treatments over time, response may depend
on treatment order.

Possible order effects include:

\[
\boxed{
\text{learning},
}
\]

\[
\boxed{
\text{fatigue},
}
\]

and

\[
\boxed{
\text{adaptation}.
}
\]

Randomizing or counterbalancing treatment order can reduce systematic
bias from these effects.

%%%%%%%%%%%%%%%%%%%%%%%%%%%%%%%%%%%%%%%%%%%%%%%%%%%%%%%%%%%%
\subsection{Carryover Effects}
\label{subsec:carryover-effects}
%%%%%%%%%%%%%%%%%%%%%%%%%%%%%%%%%%%%%%%%%%%%%%%%%%%%%%%%%%%%

A carryover effect occurs when an earlier treatment continues to affect
responses during a later treatment period.

This is especially important in crossover experiments.

A washout period may be used to reduce carryover when scientifically
appropriate.

%%%%%%%%%%%%%%%%%%%%%%%%%%%%%%%%%%%%%%%%%%%%%%%%%%%%%%%%%%%%
\subsection{Crossover Designs}
\label{subsec:crossover-designs}
%%%%%%%%%%%%%%%%%%%%%%%%%%%%%%%%%%%%%%%%%%%%%%%%%%%%%%%%%%%%

In a crossover design, experimental units receive multiple treatments in
different periods.

For two treatments \(A\) and \(B\), common sequences are

\[
\boxed{
AB
}
\]

and

\[
\boxed{
BA.
}
\]

Each participant may therefore serve as their own comparison.

%%%%%%%%%%%%%%%%%%%%%%%%%%%%%%%%%%%%%%%%%%%%%%%%%%%%%%%%%%%%
\subsection{Advantages of Crossover Designs}
\label{subsec:crossover-advantages}
%%%%%%%%%%%%%%%%%%%%%%%%%%%%%%%%%%%%%%%%%%%%%%%%%%%%%%%%%%%%

Crossover designs can reduce between-subject variability because each
unit receives multiple treatments.

They are most useful when:

\[
\boxed{
\text{conditions are stable},
}
\]

and

\[
\boxed{
\text{treatment effects are temporary}.
}
\]

They are inappropriate when treatments have long-lasting or irreversible
effects.

%%%%%%%%%%%%%%%%%%%%%%%%%%%%%%%%%%%%%%%%%%%%%%%%%%%%%%%%%%%%
\subsection{Cluster-Randomized Experiments}
\label{subsec:cluster-randomized-experiments}
%%%%%%%%%%%%%%%%%%%%%%%%%%%%%%%%%%%%%%%%%%%%%%%%%%%%%%%%%%%%

Sometimes treatments are assigned to groups rather than individuals.

Examples include randomizing:

\[
\boxed{
\text{schools},
}
\]

\[
\boxed{
\text{villages},
}
\]

\[
\boxed{
\text{hospitals},
}
\]

or

\[
\boxed{
\text{worksites}.
}
\]

Individuals within the same cluster are generally correlated.

%%%%%%%%%%%%%%%%%%%%%%%%%%%%%%%%%%%%%%%%%%%%%%%%%%%%%%%%%%%%
\subsection{Effective Sample Size in Cluster Designs}
\label{subsec:effective-sample-size-cluster}
%%%%%%%%%%%%%%%%%%%%%%%%%%%%%%%%%%%%%%%%%%%%%%%%%%%%%%%%%%%%

Suppose each cluster contains approximately \(m\) observations and the
intraclass correlation is \(\rho\).

A common approximate design effect is

\[
\boxed{
DEFF
=
1+(m-1)\rho.
}
\]

Then an approximate effective sample size is

\[
\boxed{
n_{\text{eff}}
\approx
\frac{
n
}{
DEFF
}.
}
\]

Positive within-cluster correlation therefore reduces the amount of
independent information.

%%%%%%%%%%%%%%%%%%%%%%%%%%%%%%%%%%%%%%%%%%%%%%%%%%%%%%%%%%%%
\subsection{Interference}
\label{subsec:experimental-interference}
%%%%%%%%%%%%%%%%%%%%%%%%%%%%%%%%%%%%%%%%%%%%%%%%%%%%%%%%%%%%

Many simple causal analyses assume one unit's treatment does not affect
another unit's outcome.

This may fail when units interact.

Examples include:

\[
\boxed{
\text{vaccination},
}
\]

\[
\boxed{
\text{social networks},
}
\]

\[
\boxed{
\text{classroom interventions},
}
\]

or

\[
\boxed{
\text{contagious behavior}.
}
\]

%%%%%%%%%%%%%%%%%%%%%%%%%%%%%%%%%%%%%%%%%%%%%%%%%%%%%%%%%%%%
\subsection{SUTVA Preview}
\label{subsec:sutva-preview}
%%%%%%%%%%%%%%%%%%%%%%%%%%%%%%%%%%%%%%%%%%%%%%%%%%%%%%%%%%%%

The stable-unit treatment-value assumption, abbreviated SUTVA, commonly
includes two ideas:

\[
\boxed{
\text{no interference between units},
}
\]

and

\[
\boxed{
\text{well-defined treatment versions}.
}
\]

Violations require more specialized causal models.

%%%%%%%%%%%%%%%%%%%%%%%%%%%%%%%%%%%%%%%%%%%%%%%%%%%%%%%%%%%%
\subsection{Blinding}
\label{subsec:blinding-design}
%%%%%%%%%%%%%%%%%%%%%%%%%%%%%%%%%%%%%%%%%%%%%%%%%%%%%%%%%%%%

Blinding reduces behavior or measurement changes caused by knowledge of
treatment assignment.

Possible blinded parties include:

\[
\boxed{
\text{participants},
}
\]

\[
\boxed{
\text{treatment providers},
}
\]

and

\[
\boxed{
\text{outcome assessors}.
}
\]

Blinding is especially valuable when responses contain subjective
components.

%%%%%%%%%%%%%%%%%%%%%%%%%%%%%%%%%%%%%%%%%%%%%%%%%%%%%%%%%%%%
\subsection{Placebo Control}
\label{subsec:placebo-control-design}
%%%%%%%%%%%%%%%%%%%%%%%%%%%%%%%%%%%%%%%%%%%%%%%%%%%%%%%%%%%%

A placebo resembles the active treatment without containing the active
intervention.

A placebo control can help distinguish treatment-specific effects from:

\[
\boxed{
\text{expectation},
}
\]

\[
\boxed{
\text{attention},
}
\]

or

\[
\boxed{
\text{study participation effects}.
}
\]

Placebo use must be scientifically and ethically appropriate.

%%%%%%%%%%%%%%%%%%%%%%%%%%%%%%%%%%%%%%%%%%%%%%%%%%%%%%%%%%%%
\subsection{Positive Control}
\label{subsec:positive-control}
%%%%%%%%%%%%%%%%%%%%%%%%%%%%%%%%%%%%%%%%%%%%%%%%%%%%%%%%%%%%

A positive control uses an established active treatment rather than an
inactive placebo.

This is useful when withholding effective treatment would be
inappropriate.

The scientific comparison may concern:

\[
\boxed{
\text{superiority},
}
\]

\[
\boxed{
\text{equivalence},
}
\]

or

\[
\boxed{
\text{noninferiority}.
}
\]

%%%%%%%%%%%%%%%%%%%%%%%%%%%%%%%%%%%%%%%%%%%%%%%%%%%%%%%%%%%%
\subsection{Negative Control Preview}
\label{subsec:negative-control}
%%%%%%%%%%%%%%%%%%%%%%%%%%%%%%%%%%%%%%%%%%%%%%%%%%%%%%%%%%%%

A negative control is designed so that no meaningful treatment effect is
expected.

Unexpected differences in the negative control can reveal:

\[
\boxed{
\text{measurement bias},
}
\]

\[
\boxed{
\text{confounding},
}
\]

or

\[
\boxed{
\text{procedural contamination}.
}
\]

Negative controls are especially useful in observational and laboratory
research.

%%%%%%%%%%%%%%%%%%%%%%%%%%%%%%%%%%%%%%%%%%%%%%%%%%%%%%%%%%%%
\subsection{Treatment Fidelity}
\label{subsec:treatment-fidelity}
%%%%%%%%%%%%%%%%%%%%%%%%%%%%%%%%%%%%%%%%%%%%%%%%%%%%%%%%%%%%

Treatment fidelity concerns whether the assigned treatment was delivered
as intended.

Problems include:

\[
\boxed{
\text{noncompliance},
}
\]

\[
\boxed{
\text{protocol deviations},
}
\]

and

\[
\boxed{
\text{cross-contamination}.
}
\]

Poor fidelity can reduce the distinction between treatment groups.

%%%%%%%%%%%%%%%%%%%%%%%%%%%%%%%%%%%%%%%%%%%%%%%%%%%%%%%%%%%%
\subsection{Noncompliance}
\label{subsec:experimental-noncompliance}
%%%%%%%%%%%%%%%%%%%%%%%%%%%%%%%%%%%%%%%%%%%%%%%%%%%%%%%%%%%%

Noncompliance occurs when an experimental unit does not receive or follow
the assigned treatment.

Examples include:

\[
\boxed{
\text{not taking assigned medication},
}
\]

or

\[
\boxed{
\text{switching treatments}.
}
\]

Randomization applies to treatment assignment, not necessarily treatment
received.

%%%%%%%%%%%%%%%%%%%%%%%%%%%%%%%%%%%%%%%%%%%%%%%%%%%%%%%%%%%%
\subsection{Intention-to-Treat Principle Preview}
\label{subsec:intention-to-treat-preview}
%%%%%%%%%%%%%%%%%%%%%%%%%%%%%%%%%%%%%%%%%%%%%%%%%%%%%%%%%%%%

The intention-to-treat principle analyzes participants according to their
original randomized assignment, regardless of later adherence.

This preserves the comparability created by randomization.

It often estimates the effect of treatment assignment or treatment
policy rather than the effect among perfectly compliant participants.

%%%%%%%%%%%%%%%%%%%%%%%%%%%%%%%%%%%%%%%%%%%%%%%%%%%%%%%%%%%%
\subsection{Per-Protocol Analysis Preview}
\label{subsec:per-protocol-preview}
%%%%%%%%%%%%%%%%%%%%%%%%%%%%%%%%%%%%%%%%%%%%%%%%%%%%%%%%%%%%

Per-protocol analysis focuses on units that followed the assigned
treatment protocol sufficiently closely.

Because adherence may be associated with outcomes, excluding
noncompliant units can break randomization-based comparability.

Thus per-protocol estimates can be biased without additional assumptions
or methods.

%%%%%%%%%%%%%%%%%%%%%%%%%%%%%%%%%%%%%%%%%%%%%%%%%%%%%%%%%%%%
\subsection{Attrition}
\label{subsec:experimental-attrition}
%%%%%%%%%%%%%%%%%%%%%%%%%%%%%%%%%%%%%%%%%%%%%%%%%%%%%%%%%%%%

Attrition occurs when experimental units leave the study or fail to
provide outcome measurements.

Attrition can reduce precision.

It can also create bias if dropout differs by treatment or depends on
potential outcomes.

Therefore attrition rates and reasons should be reported.

%%%%%%%%%%%%%%%%%%%%%%%%%%%%%%%%%%%%%%%%%%%%%%%%%%%%%%%%%%%%
\subsection{Missing Outcomes}
\label{subsec:missing-outcomes-design}
%%%%%%%%%%%%%%%%%%%%%%%%%%%%%%%%%%%%%%%%%%%%%%%%%%%%%%%%%%%%

Missing outcome data may arise from:

\[
\boxed{
\text{dropout},
}
\]

\[
\boxed{
\text{equipment failure},
}
\]

\[
\boxed{
\text{missed visits},
}
\]

or

\[
\boxed{
\text{data loss}.
}
\]

The statistical consequences depend on the missingness mechanism.

Missing-data methods are developed more fully in computational and
multivariate statistics.

%%%%%%%%%%%%%%%%%%%%%%%%%%%%%%%%%%%%%%%%%%%%%%%%%%%%%%%%%%%%
\subsection{Experimental Power}
\label{subsec:experimental-power}
%%%%%%%%%%%%%%%%%%%%%%%%%%%%%%%%%%%%%%%%%%%%%%%%%%%%%%%%%%%%

Power is

\[
\boxed{
P(\text{reject }H_0\mid H_A\text{ is true}).
}
\]

Experimental design strongly influences power.

Power typically increases with:

\[
\boxed{
\text{larger sample size},
}
\]

\[
\boxed{
\text{larger treatment effects},
}
\]

\[
\boxed{
\text{smaller residual variability},
}
\]

and

\[
\boxed{
\text{more efficient designs}.
}
\]

%%%%%%%%%%%%%%%%%%%%%%%%%%%%%%%%%%%%%%%%%%%%%%%%%%%%%%%%%%%%
\subsection{Sample Size Planning}
\label{subsec:sample-size-planning-design}
%%%%%%%%%%%%%%%%%%%%%%%%%%%%%%%%%%%%%%%%%%%%%%%%%%%%%%%%%%%%

Before collecting data, sample size should be linked to:

\[
\boxed{
\text{target effect size},
}
\]

\[
\boxed{
\text{desired power},
}
\]

\[
\boxed{
\text{significance level},
}
\]

\[
\boxed{
\text{expected variance},
}
\]

and

\[
\boxed{
\text{design structure}.
}
\]

Thus sample size is a design decision, not merely an arbitrary target.

%%%%%%%%%%%%%%%%%%%%%%%%%%%%%%%%%%%%%%%%%%%%%%%%%%%%%%%%%%%%
\subsection{Minimum Detectable Effect}
\label{subsec:minimum-detectable-effect}
%%%%%%%%%%%%%%%%%%%%%%%%%%%%%%%%%%%%%%%%%%%%%%%%%%%%%%%%%%%%

For fixed sample size and significance level, one may solve for the
smallest effect that can be detected with specified power.

This is the minimum detectable effect.

Conceptually,

\[
\boxed{
\text{smaller detectable effects}
\Longleftrightarrow
\text{larger or more efficient experiments}.
}
\]

%%%%%%%%%%%%%%%%%%%%%%%%%%%%%%%%%%%%%%%%%%%%%%%%%%%%%%%%%%%%
\subsection{Precision Versus Power}
\label{subsec:precision-versus-power-design}
%%%%%%%%%%%%%%%%%%%%%%%%%%%%%%%%%%%%%%%%%%%%%%%%%%%%%%%%%%%%

Power concerns the probability of detecting a specified effect.

Precision concerns the width of confidence intervals or variance of
estimators.

Both generally improve as sample size increases.

However, a design should be planned around the scientific goal rather
than around significance alone.

%%%%%%%%%%%%%%%%%%%%%%%%%%%%%%%%%%%%%%%%%%%%%%%%%%%%%%%%%%%%
\subsection{Design Efficiency}
\label{subsec:design-efficiency}
%%%%%%%%%%%%%%%%%%%%%%%%%%%%%%%%%%%%%%%%%%%%%%%%%%%%%%%%%%%%

Two designs may estimate the same treatment effect with different
variances.

Relative efficiency can be defined schematically as

\[
\boxed{
\operatorname{Eff}(D_1,D_2)
=
\frac{
\operatorname{Var}_{D_2}(\widehat{\tau})
}{
\operatorname{Var}_{D_1}(\widehat{\tau})
}.
}
\]

If the ratio exceeds \(1\), design \(D_1\) is more precise for the chosen
effect.

%%%%%%%%%%%%%%%%%%%%%%%%%%%%%%%%%%%%%%%%%%%%%%%%%%%%%%%%%%%%
\subsection{Covariate Adjustment}
\label{subsec:covariate-adjustment-design}
%%%%%%%%%%%%%%%%%%%%%%%%%%%%%%%%%%%%%%%%%%%%%%%%%%%%%%%%%%%%

Pre-treatment covariates related to the response can improve precision.

A regression-adjusted model might be

\[
\boxed{
Y
=
\beta_0
+
\tau T
+
\beta_1X
+
\varepsilon,
}
\]

where:

\[
T
=
\text{randomized treatment indicator},
\]

and

\[
X
=
\text{baseline covariate}.
\]

Randomization identifies the treatment comparison, while covariate
adjustment can reduce residual variance.

%%%%%%%%%%%%%%%%%%%%%%%%%%%%%%%%%%%%%%%%%%%%%%%%%%%%%%%%%%%%
\subsection{Stratified Randomization}
\label{subsec:stratified-randomization}
%%%%%%%%%%%%%%%%%%%%%%%%%%%%%%%%%%%%%%%%%%%%%%%%%%%%%%%%%%%%

Stratified randomization randomizes treatment separately within levels of
important baseline variables.

For example, participants might be stratified by:

\[
\boxed{
\text{site},
}
\]

\[
\boxed{
\text{age group},
}
\]

or

\[
\boxed{
\text{baseline severity}.
}
\]

This helps prevent severe treatment imbalance within important subgroups.

%%%%%%%%%%%%%%%%%%%%%%%%%%%%%%%%%%%%%%%%%%%%%%%%%%%%%%%%%%%%
\subsection{Restricted Randomization}
\label{subsec:restricted-randomization}
%%%%%%%%%%%%%%%%%%%%%%%%%%%%%%%%%%%%%%%%%%%%%%%%%%%%%%%%%%%%

Pure randomization can occasionally produce substantial treatment-group
imbalance.

Restricted randomization limits allowable assignments.

Examples include:

\[
\boxed{
\text{blocked randomization},
}
\]

\[
\boxed{
\text{stratified randomization},
}
\]

and

\[
\boxed{
\text{rerandomization}.
}
\]

The analysis should reflect the actual randomization design.

%%%%%%%%%%%%%%%%%%%%%%%%%%%%%%%%%%%%%%%%%%%%%%%%%%%%%%%%%%%%
\subsection{Rerandomization Preview}
\label{subsec:rerandomization-preview}
%%%%%%%%%%%%%%%%%%%%%%%%%%%%%%%%%%%%%%%%%%%%%%%%%%%%%%%%%%%%

In rerandomization, a random assignment is accepted only if it satisfies
a prespecified balance criterion.

If not, another random assignment is generated.

This can improve covariate balance while preserving randomization.

The acceptance rule must be defined before examining outcomes.

%%%%%%%%%%%%%%%%%%%%%%%%%%%%%%%%%%%%%%%%%%%%%%%%%%%%%%%%%%%%
\subsection{Cluster Allocation}
\label{subsec:cluster-allocation}
%%%%%%%%%%%%%%%%%%%%%%%%%%%%%%%%%%%%%%%%%%%%%%%%%%%%%%%%%%%%

In cluster-randomized studies, the number of clusters often matters more
for precision than the raw number of individuals.

For example, comparing:

\[
\boxed{
4\text{ clusters with }100\text{ people each}
}
\]

with

\[
\boxed{
40\text{ clusters with }10\text{ people each}
}
\]

may strongly favor the second design statistically when intracluster
correlation is positive.

Independent information comes primarily from independent randomized
clusters.

%%%%%%%%%%%%%%%%%%%%%%%%%%%%%%%%%%%%%%%%%%%%%%%%%%%%%%%%%%%%
\subsection{Blocking Versus Stratification}
\label{subsec:blocking-versus-stratification}
%%%%%%%%%%%%%%%%%%%%%%%%%%%%%%%%%%%%%%%%%%%%%%%%%%%%%%%%%%%%

In sampling, stratification divides the population into subgroups before
sampling.

In experiments, blocking divides experimental units into similar groups
before random treatment assignment.

The mathematical ideas are closely related:

\[
\boxed{
\text{control known heterogeneity before randomization}.
}
\]

%%%%%%%%%%%%%%%%%%%%%%%%%%%%%%%%%%%%%%%%%%%%%%%%%%%%%%%%%%%%
\subsection{Blocking Versus Covariate Adjustment}
\label{subsec:blocking-versus-covariate-adjustment}
%%%%%%%%%%%%%%%%%%%%%%%%%%%%%%%%%%%%%%%%%%%%%%%%%%%%%%%%%%%%

Blocking uses a variable during the design stage.

Covariate adjustment uses a variable during the analysis stage.

Both can improve precision.

When possible, using an important variable in both design and analysis
can be effective, provided the procedure is properly specified.

%%%%%%%%%%%%%%%%%%%%%%%%%%%%%%%%%%%%%%%%%%%%%%%%%%%%%%%%%%%%
\subsection{Randomization Inference}
\label{subsec:randomization-inference-design}
%%%%%%%%%%%%%%%%%%%%%%%%%%%%%%%%%%%%%%%%%%%%%%%%%%%%%%%%%%%%

Randomization itself induces a probability distribution over treatment
assignments.

Under a sharp null hypothesis of no treatment effect, outcomes can be
treated as fixed while treatment labels vary according to the design.

This produces a randomization distribution for a test statistic.

Thus statistical inference can sometimes be based directly on the
experimental design rather than on a population-distribution model.

%%%%%%%%%%%%%%%%%%%%%%%%%%%%%%%%%%%%%%%%%%%%%%%%%%%%%%%%%%%%
\subsection{Sharp Null Hypothesis}
\label{subsec:sharp-null-hypothesis}
%%%%%%%%%%%%%%%%%%%%%%%%%%%%%%%%%%%%%%%%%%%%%%%%%%%%%%%%%%%%

A sharp null specifies every missing potential outcome.

A common example is:

\[
\boxed{
Y_i(1)=Y_i(0)
\quad
\text{for every unit }i.
}
\]

This states that treatment has no effect on any unit.

Such a null permits exact randomization testing because all potential
outcomes are determined under the hypothesis.

%%%%%%%%%%%%%%%%%%%%%%%%%%%%%%%%%%%%%%%%%%%%%%%%%%%%%%%%%%%%
\subsection{Average Treatment Effect}
\label{subsec:average-treatment-effect-design}
%%%%%%%%%%%%%%%%%%%%%%%%%%%%%%%%%%%%%%%%%%%%%%%%%%%%%%%%%%%%

In potential-outcomes notation, define

\[
Y_i(1)
\]

as unit \(i\)'s response under treatment and

\[
Y_i(0)
\]

under control.

The individual treatment effect is

\[
\boxed{
\tau_i
=
Y_i(1)-Y_i(0).
}
\]

The finite-population average treatment effect is

\[
\boxed{
\tau
=
\frac1N
\sum_{i=1}^{N}
\left[
Y_i(1)-Y_i(0)
\right].
}
\]

%%%%%%%%%%%%%%%%%%%%%%%%%%%%%%%%%%%%%%%%%%%%%%%%%%%%%%%%%%%%
\subsection{Difference-in-Means Estimator}
\label{subsec:difference-in-means-estimator-design}
%%%%%%%%%%%%%%%%%%%%%%%%%%%%%%%%%%%%%%%%%%%%%%%%%%%%%%%%%%%%

For a randomized treatment-control experiment, a natural estimator is

\[
\boxed{
\widehat{\tau}
=
\overline{Y}_T
-
\overline{Y}_C.
}
\]

Under complete randomization, this estimator is unbiased for the
finite-population average treatment effect.

Randomization therefore gives a direct causal interpretation to the
difference in sample means.

%%%%%%%%%%%%%%%%%%%%%%%%%%%%%%%%%%%%%%%%%%%%%%%%%%%%%%%%%%%%
\subsection{Potential Outcomes and Missing Data}
\label{subsec:potential-outcomes-missing}
%%%%%%%%%%%%%%%%%%%%%%%%%%%%%%%%%%%%%%%%%%%%%%%%%%%%%%%%%%%%

For each unit, only one potential outcome is observed.

If unit \(i\) receives treatment,

\[
Y_i(1)
\]

is observed but

\[
Y_i(0)
\]

is missing.

If unit \(i\) receives control, the reverse occurs.

Thus causal inference contains a fundamental missing-data problem.

Randomization allows treatment groups to represent each other's
unobserved counterfactual outcomes in aggregate.

%%%%%%%%%%%%%%%%%%%%%%%%%%%%%%%%%%%%%%%%%%%%%%%%%%%%%%%%%%%%
\subsection{Sequential Experiments Preview}
\label{subsec:sequential-experiments-preview}
%%%%%%%%%%%%%%%%%%%%%%%%%%%%%%%%%%%%%%%%%%%%%%%%%%%%%%%%%%%%

In a sequential experiment, data are examined during collection and the
study may adapt or stop according to a prespecified rule.

Examples include:

\[
\boxed{
\text{group-sequential trials},
}
\]

\[
\boxed{
\text{sequential probability ratio tests},
}
\]

and

\[
\boxed{
\text{adaptive allocation}.
}
\]

Ordinary fixed-sample significance thresholds cannot generally be reused
without adjustment.

%%%%%%%%%%%%%%%%%%%%%%%%%%%%%%%%%%%%%%%%%%%%%%%%%%%%%%%%%%%%
\subsection{Adaptive Designs Preview}
\label{subsec:adaptive-designs-preview}
%%%%%%%%%%%%%%%%%%%%%%%%%%%%%%%%%%%%%%%%%%%%%%%%%%%%%%%%%%%%

An adaptive design allows certain features of the experiment to change
using accumulated information.

Possible adaptations include:

\[
\boxed{
\text{sample-size modification},
}
\]

\[
\boxed{
\text{dropping ineffective treatments},
}
\]

\[
\boxed{
\text{changing allocation probabilities},
}
\]

or

\[
\boxed{
\text{stopping early}.
}
\]

Valid adaptive designs specify these rules in advance and preserve
appropriate error control.

%%%%%%%%%%%%%%%%%%%%%%%%%%%%%%%%%%%%%%%%%%%%%%%%%%%%%%%%%%%%
\subsection{Response-Adaptive Randomization Preview}
\label{subsec:response-adaptive-randomization}
%%%%%%%%%%%%%%%%%%%%%%%%%%%%%%%%%%%%%%%%%%%%%%%%%%%%%%%%%%%%

Response-adaptive randomization changes future treatment probabilities
based on outcomes already observed.

The goal may be to assign more participants to better-performing
treatments.

However, adaptive allocation changes the sampling structure and can make
analysis more complex.

%%%%%%%%%%%%%%%%%%%%%%%%%%%%%%%%%%%%%%%%%%%%%%%%%%%%%%%%%%%%
\subsection{Optimal Experimental Design Preview}
\label{subsec:optimal-experimental-design-preview}
%%%%%%%%%%%%%%%%%%%%%%%%%%%%%%%%%%%%%%%%%%%%%%%%%%%%%%%%%%%%

Experimental design can itself be formulated as an optimization problem.

Suppose the regression model is

\[
\mathbf{y}
=
X\boldsymbol{\beta}
+
\boldsymbol{\varepsilon}.
\]

Then

\[
\operatorname{Var}
(
\widehat{\boldsymbol{\beta}}
)
=
\sigma^2
(X^TX)^{-1}.
\]

A design may therefore choose \(X\) to optimize a function of

\[
(X^TX)^{-1}.
\]

%%%%%%%%%%%%%%%%%%%%%%%%%%%%%%%%%%%%%%%%%%%%%%%%%%%%%%%%%%%%
\subsection{D-Optimality Preview}
\label{subsec:d-optimality-preview}
%%%%%%%%%%%%%%%%%%%%%%%%%%%%%%%%%%%%%%%%%%%%%%%%%%%%%%%%%%%%

A D-optimal design seeks to maximize

\[
\boxed{
\det(X^TX).
}
\]

Equivalently, it minimizes the determinant of the covariance matrix up to
a multiplicative constant.

Geometrically, this tends to minimize the volume of a joint confidence
ellipsoid for the coefficients.

%%%%%%%%%%%%%%%%%%%%%%%%%%%%%%%%%%%%%%%%%%%%%%%%%%%%%%%%%%%%
\subsection{A-Optimality Preview}
\label{subsec:a-optimality-preview}
%%%%%%%%%%%%%%%%%%%%%%%%%%%%%%%%%%%%%%%%%%%%%%%%%%%%%%%%%%%%

An A-optimal design seeks to minimize

\[
\boxed{
\operatorname{tr}
\left[
(X^TX)^{-1}
\right].
}
\]

This approximately minimizes the average coefficient variance under the
classical linear model.

Different optimality criteria emphasize different inferential goals.

%%%%%%%%%%%%%%%%%%%%%%%%%%%%%%%%%%%%%%%%%%%%%%%%%%%%%%%%%%%%
\subsection{Response Surface Designs Preview}
\label{subsec:response-surface-designs-preview}
%%%%%%%%%%%%%%%%%%%%%%%%%%%%%%%%%%%%%%%%%%%%%%%%%%%%%%%%%%%%

When quantitative factors influence a response smoothly, one may fit a
second-order model such as

\[
\boxed{
Y
=
\beta_0
+
\sum_i
\beta_iX_i
+
\sum_i
\beta_{ii}X_i^2
+
\sum_{i<j}
\beta_{ij}X_iX_j
+
\varepsilon.
}
\]

Designs for estimating such models efficiently are studied in response
surface methodology.

%%%%%%%%%%%%%%%%%%%%%%%%%%%%%%%%%%%%%%%%%%%%%%%%%%%%%%%%%%%%
\subsection{Central Composite Design Preview}
\label{subsec:central-composite-design-preview}
%%%%%%%%%%%%%%%%%%%%%%%%%%%%%%%%%%%%%%%%%%%%%%%%%%%%%%%%%%%%

A central composite design combines:

\[
\boxed{
\text{factorial points},
}
\]

\[
\boxed{
\text{axial points},
}
\]

and often

\[
\boxed{
\text{center points}.
}
\]

It is commonly used to estimate quadratic response surfaces without
requiring a full grid of factor combinations.

%%%%%%%%%%%%%%%%%%%%%%%%%%%%%%%%%%%%%%%%%%%%%%%%%%%%%%%%%%%%
\subsection{Center Points}
\label{subsec:center-points}
%%%%%%%%%%%%%%%%%%%%%%%%%%%%%%%%%%%%%%%%%%%%%%%%%%%%%%%%%%%%

In quantitative-factor experiments, repeated center points can provide:

\[
\boxed{
\text{pure-error estimation},
}
\]

and

\[
\boxed{
\text{information about curvature}.
}
\]

If the average center response differs systematically from the fitted
first-order model, curvature may be present.

%%%%%%%%%%%%%%%%%%%%%%%%%%%%%%%%%%%%%%%%%%%%%%%%%%%%%%%%%%%%
\subsection{Screening Experiments}
\label{subsec:screening-experiments}
%%%%%%%%%%%%%%%%%%%%%%%%%%%%%%%%%%%%%%%%%%%%%%%%%%%%%%%%%%%%

A screening experiment identifies which factors among many candidates
have important effects.

The goal is often not detailed optimization but dimensional reduction:

\[
\boxed{
\text{many candidate factors}
\longrightarrow
\text{small set of active factors}.
}
\]

Fractional factorial and related designs are commonly used.

%%%%%%%%%%%%%%%%%%%%%%%%%%%%%%%%%%%%%%%%%%%%%%%%%%%%%%%%%%%%
\subsection{Confirmation Experiments}
\label{subsec:confirmation-experiments}
%%%%%%%%%%%%%%%%%%%%%%%%%%%%%%%%%%%%%%%%%%%%%%%%%%%%%%%%%%%%

After identifying a promising treatment combination or model, a separate
confirmation experiment can test whether the predicted improvement is
reproducible.

This protects against selecting a configuration that appears favorable
only because of random noise or model overfitting.

%%%%%%%%%%%%%%%%%%%%%%%%%%%%%%%%%%%%%%%%%%%%%%%%%%%%%%%%%%%%
\subsection{Internal Validity}
\label{subsec:internal-validity-experiment}
%%%%%%%%%%%%%%%%%%%%%%%%%%%%%%%%%%%%%%%%%%%%%%%%%%%%%%%%%%%%

Internal validity concerns whether treatment comparisons within the study
support the intended causal interpretation.

Threats include:

\[
\boxed{
\text{failed randomization},
}
\]

\[
\boxed{
\text{differential attrition},
}
\]

\[
\boxed{
\text{noncompliance},
}
\]

\[
\boxed{
\text{measurement bias},
}
\]

and

\[
\boxed{
\text{interference}.
}
\]

%%%%%%%%%%%%%%%%%%%%%%%%%%%%%%%%%%%%%%%%%%%%%%%%%%%%%%%%%%%%
\subsection{External Validity}
\label{subsec:external-validity-experiment}
%%%%%%%%%%%%%%%%%%%%%%%%%%%%%%%%%%%%%%%%%%%%%%%%%%%%%%%%%%%%

External validity concerns whether experimental conclusions generalize
to:

\[
\boxed{
\text{other populations},
}
\]

\[
\boxed{
\text{other settings},
}
\]

\[
\boxed{
\text{other times},
}
\]

or

\[
\boxed{
\text{other treatment implementations}.
}
\]

Random assignment does not automatically guarantee external validity.

%%%%%%%%%%%%%%%%%%%%%%%%%%%%%%%%%%%%%%%%%%%%%%%%%%%%%%%%%%%%
\subsection{Ecological Validity}
\label{subsec:ecological-validity}
%%%%%%%%%%%%%%%%%%%%%%%%%%%%%%%%%%%%%%%%%%%%%%%%%%%%%%%%%%%%

An experiment conducted under highly controlled conditions may differ
substantially from real-world settings.

Ecological validity concerns whether the experimental environment
resembles the context where conclusions will be applied.

Control and realism often involve tradeoffs.

%%%%%%%%%%%%%%%%%%%%%%%%%%%%%%%%%%%%%%%%%%%%%%%%%%%%%%%%%%%%
\subsection{Experimental Design Workflow}
\label{subsec:experimental-design-workflow}
%%%%%%%%%%%%%%%%%%%%%%%%%%%%%%%%%%%%%%%%%%%%%%%%%%%%%%%%%%%%

A systematic workflow is:

\begin{enumerate}
    \item define the scientific question;
    \item define the target population;
    \item define the response variable;
    \item identify experimental units;
    \item identify treatment factors and levels;
    \item identify nuisance variables;
    \item determine whether blocking is useful;
    \item determine the randomization unit;
    \item select the treatment-allocation structure;
    \item determine replication;
    \item perform power or precision calculations;
    \item define the randomization procedure;
    \item specify blinding when possible;
    \item define data-collection protocols;
    \item define exclusion and missing-data procedures;
    \item specify the primary analysis before observing outcomes;
    \item conduct the experiment;
    \item verify treatment implementation;
    \item analyze according to the actual design;
    \item report effect estimates, uncertainty, and design limitations.
\end{enumerate}

%%%%%%%%%%%%%%%%%%%%%%%%%%%%%%%%%%%%%%%%%%%%%%%%%%%%%%%%%%%%
\subsection{Choosing a Design}
\label{subsec:choosing-experimental-design}
%%%%%%%%%%%%%%%%%%%%%%%%%%%%%%%%%%%%%%%%%%%%%%%%%%%%%%%%%%%%

Use a completely randomized design when:

\[
\boxed{
\text{units are reasonably homogeneous}
}
\]

and unrestricted randomization is practical.

Use blocking when:

\[
\boxed{
\text{known nuisance variables explain important response variation}.
}
\]

Use factorial designs when:

\[
\boxed{
\text{several factors and interactions are scientifically important}.
}
\]

Use split-plot or cluster designs when:

\[
\boxed{
\text{randomization must occur at multiple physical levels}.
}
\]

%%%%%%%%%%%%%%%%%%%%%%%%%%%%%%%%%%%%%%%%%%%%%%%%%%%%%%%%%%%%
\subsection{Common Errors}
\label{subsec:experimental-design-common-errors}
%%%%%%%%%%%%%%%%%%%%%%%%%%%%%%%%%%%%%%%%%%%%%%%%%%%%%%%%%%%%

\subsubsection{Confusing Observational Units with Experimental Units}

If treatment is randomized by classroom, then the classroom is the
experimental unit even if many students are measured.

%%%%%%%%%%%%%%%%%%%%%%%%%%%%%%%%%%%%%%%%%%%%%%%%%%%%%%%%%%%%
\subsubsection{Treating Repeated Measurements as Independent Replication}

Repeated measurements from one unit do not create multiple independent
treatment assignments.

%%%%%%%%%%%%%%%%%%%%%%%%%%%%%%%%%%%%%%%%%%%%%%%%%%%%%%%%%%%%
\subsubsection{Ignoring Pseudoreplication}

Analyzing subsamples as independent replicates can dramatically
underestimate uncertainty.

%%%%%%%%%%%%%%%%%%%%%%%%%%%%%%%%%%%%%%%%%%%%%%%%%%%%%%%%%%%%
\subsubsection{Ignoring Randomization Level}

The correct analysis must match the level at which treatment was actually
randomized.

%%%%%%%%%%%%%%%%%%%%%%%%%%%%%%%%%%%%%%%%%%%%%%%%%%%%%%%%%%%%
\subsubsection{Blocking on Irrelevant Variables}

Blocking only improves precision when the blocking variable meaningfully
explains response variation.

Too many unnecessary blocks can complicate the design.

%%%%%%%%%%%%%%%%%%%%%%%%%%%%%%%%%%%%%%%%%%%%%%%%%%%%%%%%%%%%
\subsubsection{Blocking on Post-Treatment Variables}

Blocking is a design-stage procedure and should use pre-treatment
information.

Variables affected by treatment cannot be used for pre-treatment
blocking.

%%%%%%%%%%%%%%%%%%%%%%%%%%%%%%%%%%%%%%%%%%%%%%%%%%%%%%%%%%%%
\subsubsection{Assuming More Factors Always Improve an Experiment}

Adding factors rapidly increases the number of treatment combinations.

A full \(2^k\) design doubles in size each time a new two-level factor is
added.

%%%%%%%%%%%%%%%%%%%%%%%%%%%%%%%%%%%%%%%%%%%%%%%%%%%%%%%%%%%%
\subsubsection{Ignoring Interactions}

Studying factors one at a time can miss situations where their effects
depend on one another.

%%%%%%%%%%%%%%%%%%%%%%%%%%%%%%%%%%%%%%%%%%%%%%%%%%%%%%%%%%%%
\subsubsection{Interpreting Main Effects Without Checking Interaction}

Strong interaction can make an averaged main effect scientifically
misleading.

%%%%%%%%%%%%%%%%%%%%%%%%%%%%%%%%%%%%%%%%%%%%%%%%%%%%%%%%%%%%
\subsubsection{Confusing Experimental and Observational Confounding}

Experimental confounding may refer to deliberate aliasing of effects in a
fractional design.

Observational confounding concerns uncontrolled variables affecting
causal interpretation.

%%%%%%%%%%%%%%%%%%%%%%%%%%%%%%%%%%%%%%%%%%%%%%%%%%%%%%%%%%%%
\subsubsection{Assuming Randomization Guarantees Perfect Balance}

Randomization balances covariates probabilistically, not exactly in every
realized experiment.

%%%%%%%%%%%%%%%%%%%%%%%%%%%%%%%%%%%%%%%%%%%%%%%%%%%%%%%%%%%%
\subsubsection{Changing the Analysis After Seeing Outcomes}

Unplanned analyses can inflate false positive rates and create selective
reporting problems.

Primary analyses should be specified before outcome inspection when
possible.

%%%%%%%%%%%%%%%%%%%%%%%%%%%%%%%%%%%%%%%%%%%%%%%%%%%%%%%%%%%%
\subsubsection{Ignoring Attrition}

If dropout differs across treatments or relates to outcomes, the original
randomized comparison may become distorted.

%%%%%%%%%%%%%%%%%%%%%%%%%%%%%%%%%%%%%%%%%%%%%%%%%%%%%%%%%%%%
\subsubsection{Ignoring Noncompliance}

Analyzing treatment received instead of treatment assigned can introduce
selection bias.

%%%%%%%%%%%%%%%%%%%%%%%%%%%%%%%%%%%%%%%%%%%%%%%%%%%%%%%%%%%%
\subsubsection{Ignoring Cluster Correlation}

Individuals in the same randomized cluster are not independent
experimental replicates.

%%%%%%%%%%%%%%%%%%%%%%%%%%%%%%%%%%%%%%%%%%%%%%%%%%%%%%%%%%%%
\subsubsection{Using a Fixed-Sample Test After Repeated Interim Looks}

Repeated testing during data collection can inflate Type I error unless
the design accounts for sequential monitoring.

%%%%%%%%%%%%%%%%%%%%%%%%%%%%%%%%%%%%%%%%%%%%%%%%%%%%%%%%%%%%
\subsubsection{Ignoring Treatment Implementation}

A statistically elegant design cannot answer the intended question if
treatments were not delivered as specified.

%%%%%%%%%%%%%%%%%%%%%%%%%%%%%%%%%%%%%%%%%%%%%%%%%%%%%%%%%%%%
\subsection{Applications}
\label{subsec:experimental-design-applications}
%%%%%%%%%%%%%%%%%%%%%%%%%%%%%%%%%%%%%%%%%%%%%%%%%%%%%%%%%%%%

\begin{applicationbox}

\textbf{Agriculture.}
Blocks, factorial designs, and split plots are used to study fertilizers,
irrigation, varieties, soil conditions, and field heterogeneity.

\medskip

\textbf{Engineering.}
Factorial and response-surface designs identify important process factors,
interactions, and optimal operating conditions.

\medskip

\textbf{Manufacturing.}
Designed experiments improve quality, reduce defects, and identify robust
process settings.

\medskip

\textbf{Medicine.}
Randomized controlled trials compare treatments, doses, interventions,
and treatment sequences.

\medskip

\textbf{Education Research.}
Students, classrooms, or schools may be randomized to programs or
teaching interventions.

\medskip

\textbf{Environmental Science.}
Treatments may be applied to plots, ponds, habitats, monitoring stations,
or ecological systems using blocked and repeated designs.

\medskip

\textbf{Psychology.}
Between-subject, within-subject, factorial, and crossover designs are
used to isolate behavioral effects.

\medskip

\textbf{Technology and Business.}
A/B tests and multivariate experiments compare website, pricing,
advertising, and product-design changes.

\end{applicationbox}

%%%%%%%%%%%%%%%%%%%%%%%%%%%%%%%%%%%%%%%%%%%%%%%%%%%%%%%%%%%%
\subsection{Exercises}
\label{subsec:experimental-design-exercises}
%%%%%%%%%%%%%%%%%%%%%%%%%%%%%%%%%%%%%%%%%%%%%%%%%%%%%%%%%%%%

\begin{exercise}[1]
Define an experimental unit.
\end{exercise}

\begin{exercise}[1]
Define a factor.
\end{exercise}

\begin{exercise}[1]
Define a treatment.
\end{exercise}

\begin{exercise}[1]
Define replication.
\end{exercise}

\begin{exercise}[1]
Define blocking.
\end{exercise}

\begin{exercise}[1]
Define randomization.
\end{exercise}

\begin{exercise}[1]
Distinguish random sampling from random assignment.
\end{exercise}

\begin{exercise}[1]
Define pseudoreplication.
\end{exercise}

\begin{exercise}[1]
Define a completely randomized design.
\end{exercise}

\begin{exercise}[1]
Define a randomized complete block design.
\end{exercise}

\begin{exercise}[1]
Define a factorial experiment.
\end{exercise}

\begin{exercise}[1]
Define interaction.
\end{exercise}

\begin{exercise}[2]
A study randomly assigns \(60\) subjects equally to three treatments.
How many independent experimental units receive each treatment?
\end{exercise}

\begin{exercise}[2]
Four classrooms are randomized to two teaching methods, and \(25\)
students are measured in each classroom. Identify the experimental and
observational units.
\end{exercise}

\begin{exercise}[2]
Explain why the previous experiment does not have \(100\) independently
randomized experimental units.
\end{exercise}

\begin{exercise}[2]
A \(2^3\) factorial design is planned. How many treatment combinations
are required?
\end{exercise}

\begin{exercise}[2]
A \(2^5\) full factorial design is planned. How many treatment
combinations are required?
\end{exercise}

\begin{exercise}[2]
A factor has levels coded \(-1\) and \(+1\). If \(A=-1\) and \(B=+1\),
find the coded value of the \(AB\) interaction.
\end{exercise}

\begin{exercise}[2]
A \(3\times4\) factorial design has \(5\) independent replicates per
treatment combination. Find the total sample size.
\end{exercise}

\begin{exercise}[2]
A cluster-randomized experiment has average cluster size \(m=20\) and
intraclass correlation \(\rho=0.05\). Compute the approximate design
effect.
\end{exercise}

\begin{exercise}[2]
Using the previous design effect, find the approximate effective sample
size if \(n=400\) individuals are observed.
\end{exercise}

\begin{exercise}[2]
Construct a simple \(3\times3\) Latin square using treatments \(A,B,C\).
\end{exercise}

\begin{exercise}[2]
Explain why a before-and-after study on the same individuals is not an
independent-groups design.
\end{exercise}

\begin{exercise}[3]
Explain how randomization protects against systematic treatment-group
differences.
\end{exercise}

\begin{exercise}[3]
Explain why randomization does not guarantee exact balance in every
experiment.
\end{exercise}

\begin{exercise}[3]
Explain mathematically why positive within-pair correlation can improve
the precision of a matched-pairs design.
\end{exercise}

\begin{exercise}[3]
Explain why blocking can reduce residual variance.
\end{exercise}

\begin{exercise}[3]
Show that a full factorial experiment with factors having
\(a_1,\ldots,a_k\) levels requires

\[
\prod_{j=1}^{k}a_j
\]

treatment combinations.
\end{exercise}

\begin{exercise}[3]
Construct the complete design table for a \(2^3\) factorial using
\(-1,+1\) coding.
\end{exercise}

\begin{exercise}[3]
Add columns for

\[
AB,\ AC,\ BC,\ ABC
\]

to the \(2^3\) design.
\end{exercise}

\begin{exercise}[3]
Show that the main-effect columns in a balanced \(2^k\) factorial are
orthogonal.
\end{exercise}

\begin{exercise}[3]
Explain why factorial designs can reveal interactions that one-factor-at-
a-time experiments cannot.
\end{exercise}

\begin{exercise}[3]
Explain the distinction between crossed and nested factors.
\end{exercise}

\begin{exercise}[3]
Give an example of a split-plot experiment and identify the whole-plot
and subplot factors.
\end{exercise}

\begin{exercise}[3]
Explain why split-plot effects use different error terms.
\end{exercise}

\begin{exercise}[3]
Explain why repeated measurements from the same person are correlated.
\end{exercise}

\begin{exercise}[3]
Explain the purpose of counterbalancing treatment order.
\end{exercise}

\begin{exercise}[3]
Explain how carryover effects can bias a crossover experiment.
\end{exercise}

\begin{exercise}[3]
Explain why cluster randomization usually requires more total
observations than individual randomization for the same power.
\end{exercise}

\begin{exercise}[3]
Explain the distinction between internal and external validity.
\end{exercise}

\begin{exercise}[3]
Explain why treatment noncompliance can complicate causal interpretation.
\end{exercise}

\begin{exercise}[3]
Explain why intention-to-treat analysis preserves randomization.
\end{exercise}

\begin{exercise}[3]
Explain why analyzing only compliant participants may create selection
bias.
\end{exercise}

\begin{exercise}[3]
Explain why attrition can threaten internal validity.
\end{exercise}

\begin{exercise}[3]
Explain the difference between experimental confounding and observational
confounding.
\end{exercise}

\begin{exercise}[3]
If a defining relation is

\[
I=ABC,
\]

derive the aliases of \(A\), \(B\), and \(C\).
\end{exercise}

\begin{exercise}[3]
Explain the difference between Resolution III, IV, and V fractional
factorial designs.
\end{exercise}

\begin{exercise}[3]
Explain why center points help detect curvature in response-surface
experiments.
\end{exercise}

\begin{exercise}[4]
Derive the unbiasedness of the difference-in-means estimator under
complete randomization using potential outcomes.
\end{exercise}

\begin{exercise}[4]
Derive the randomization variance of the difference-in-means estimator in
a completely randomized experiment.
\end{exercise}

\begin{exercise}[4]
Study Fisher randomization tests under the sharp null hypothesis.
\end{exercise}

\begin{exercise}[4]
Study Neymanian repeated-sampling inference for randomized experiments.
\end{exercise}

\begin{exercise}[4]
Derive the variance reduction obtained by blocking under a suitable
additive model.
\end{exercise}

\begin{exercise}[4]
Study optimal allocation between two treatment groups when costs or
variances differ.
\end{exercise}

\begin{exercise}[4]
Derive factorial effect estimates for a \(2^k\) design using contrast
coefficients.
\end{exercise}

\begin{exercise}[4]
Study alias structures in regular fractional factorial designs.
\end{exercise}

\begin{exercise}[4]
Construct a Resolution IV fractional factorial design.
\end{exercise}

\begin{exercise}[4]
Investigate Plackett--Burman screening designs.
\end{exercise}

\begin{exercise}[4]
Study response-surface methodology and central composite designs.
\end{exercise}

\begin{exercise}[4]
Investigate Box--Behnken designs.
\end{exercise}

\begin{exercise}[4]
Study D-optimal and A-optimal experimental designs.
\end{exercise}

\begin{exercise}[4]
Investigate rerandomization using Mahalanobis-distance balance criteria.
\end{exercise}

\begin{exercise}[4]
Study covariate-adaptive randomization.
\end{exercise}

\begin{exercise}[4]
Investigate group-sequential experiments.
\end{exercise}

\begin{exercise}[4]
Study response-adaptive randomization.
\end{exercise}

\begin{exercise}[4]
Investigate cluster-randomized trial power and sample-size formulas.
\end{exercise}

\begin{exercise}[4]
Study interference and network experiments.
\end{exercise}

\begin{exercise}[4]
Investigate principal stratification and causal effects under
noncompliance.
\end{exercise}

%%%%%%%%%%%%%%%%%%%%%%%%%%%%%%%%%%%%%%%%%%%%%%%%%%%%%%%%%%%%
\subsection{Conceptual Summary}
\label{subsec:experimental-design-summary}
%%%%%%%%%%%%%%%%%%%%%%%%%%%%%%%%%%%%%%%%%%%%%%%%%%%%%%%%%%%%

Experimental design determines how treatments are assigned and therefore
what causal and statistical conclusions can be supported.

The basic elements are

\[
\boxed{
\text{experimental units},
}
\]

\[
\boxed{
\text{treatments},
}
\]

\[
\boxed{
\text{factors},
}
\]

and

\[
\boxed{
\text{responses}.
}
\]

Three classical principles are

\[
\boxed{
\text{randomization},
}
\]

\[
\boxed{
\text{replication},
}
\]

and

\[
\boxed{
\text{blocking}.
}
\]

Randomization supports causal inference.

Replication estimates experimental variability.

Blocking removes known nuisance variation.

A completely randomized design assigns treatments directly at random.

A randomized complete block design randomizes treatments within blocks.

Factorial experiments study several factors simultaneously and permit
estimation of interactions.

For \(k\) two-level factors, a full factorial requires

\[
\boxed{
2^k
}
\]

treatment combinations.

Fractional factorial designs reduce experimental cost by deliberately
aliasing some effects.

Cluster, repeated-measures, split-plot, crossover, and hierarchical
experiments require analysis that respects the actual randomization and
dependence structure.

Thus experimental design combines

\[
\boxed{
\text{causal reasoning}
+
\text{probability}
+
\text{randomization}
+
\text{ANOVA}
+
\text{regression}
+
\text{optimization}.
}
\]

%%%%%%%%%%%%%%%%%%%%%%%%%%%%%%%%%%%%%%%%%%%%%%%%%%%%%%%%%%%%
\subsection{Section Connections}
\label{subsec:experimental-design-connections}
%%%%%%%%%%%%%%%%%%%%%%%%%%%%%%%%%%%%%%%%%%%%%%%%%%%%%%%%%%%%

\chapterconnection{
Experimental Design provides the data-generation framework underlying
the ANOVA methods of Section~\ref{sec:analysis-of-variance}. Completely
randomized experiments lead naturally to one-way ANOVA, blocked designs
lead to block models, and factorial experiments produce main effects and
interactions within the general linear-model framework of
Section~\ref{sec:regression}. Random assignment also provides the
foundation for causal interpretation of treatment effects. The next
section, Categorical Data Analysis, shifts attention from continuous
responses to categorical outcomes and develops contingency-table,
odds-ratio, logistic, and log-linear methods. Experimental-design ideas
reappear throughout statistical learning, causal inference, optimization,
and applied mathematical modeling.
}

%%%%%%%%%%%%%%%%%%%%%%%%%%%%%%%%%%%%%%%%%%%%%%%%%%%%%%%%%%%%
% SECTION SOURCE TRACKING
%%%%%%%%%%%%%%%%%%%%%%%%%%%%%%%%%%%%%%%%%%%%%%%%%%%%%%%%%%%%

%------------------------------------------------------------
% 8.9 Experimental Design
%------------------------------------------------------------
%
% Source basis:
%   - Standard design of experiments.
%   - Standard randomized experimental design.
%   - Standard factorial, block, and response-surface concepts.
%   - Standard potential-outcomes and randomization-inference ideas.
%
% Used for:
%   - experimental units;
%   - observational units;
%   - treatments;
%   - factors and factor levels;
%   - response variables;
%   - randomization;
%   - replication;
%   - blocking;
%   - pseudoreplication;
%   - nuisance variables;
%   - completely randomized designs;
%   - balanced and unequal allocation;
%   - randomized complete block designs;
%   - matched-pairs designs;
%   - Latin square designs;
%   - factorial experiments;
%   - main effects;
%   - interactions;
%   - two-level factorials;
%   - coded factors;
%   - factorial interaction columns;
%   - hierarchy and effect sparsity;
%   - fractional factorial designs;
%   - aliasing;
%   - defining relations;
%   - design resolution;
%   - experimental confounding;
%   - orthogonality;
%   - restricted randomization;
%   - split-plot designs;
%   - nested and crossed factors;
%   - fixed and random factors;
%   - repeated-measures designs;
%   - order and carryover effects;
%   - crossover designs;
%   - cluster-randomized experiments;
%   - design effects;
%   - interference;
%   - SUTVA preview;
%   - blinding;
%   - placebo and active controls;
%   - negative controls;
%   - treatment fidelity;
%   - noncompliance;
%   - intention-to-treat;
%   - per-protocol analysis;
%   - attrition;
%   - power and sample-size planning;
%   - minimum detectable effects;
%   - design efficiency;
%   - covariate adjustment;
%   - stratified and restricted randomization;
%   - rerandomization;
%   - randomization inference;
%   - potential outcomes;
%   - average treatment effects;
%   - sequential and adaptive design previews;
%   - optimal design;
%   - D- and A-optimality;
%   - response-surface methodology;
%   - central composite designs;
%   - screening and confirmation experiments;
%   - internal and external validity;
%   - applications and exercises.
%
% Notes:
%   - Categorical Data Analysis is developed next in Section 8.10.
%   - ANOVA in Section 8.8 supplies much of the classical analysis for
%     randomized and factorial experiments.
%   - Regression in Section 8.7 provides the general linear-model
%     representation of experimental effects and interactions.
%   - Optimization and statistical learning later expand optimal-design,
%     adaptive-design, and predictive-design concepts.
%
%%%%%%%%%%%%%%%%%%%%%%%%%%%%%%%%%%%%%%%%%%%%%%%%%%%%%%%%%%%%